\section{Compact IBM Hardware Validation}
\label{app:hardware-validation}

E9 asks a narrow question: can a readout fitted from exact features retain useful
predictions when the same eight inputs are represented by finite-shot,
stochastic-trajectory features from an IBM QPU?  The study is not designed to
establish hardware-noise robustness across devices, quantum advantage, or
scalable deployment.

\subsection{Program, data, and frozen readout}
\label{app:hardware-program}

The hardware-compatible program uses a medium-persistence \(d=2\) VARMA process,
window length \(w=4\), stride \(s=100\), one two-qubit reservoir, ring/chain
edge \((0,1)\), \(\lambda=0.8\), the \(\tanh\) angle map, and
\[
 \cO_{\mathrm{HW}}
 =\{X_0,Y_0,Z_0,X_1,Y_1,Z_1,Z_0Z_1\}.
\]
The stored observable labels are
\texttt{IX,IY,IZ,XI,YI,ZI,ZZ} in the package's qubit-order convention.  The two
tasks are delayed recall at \(\tau=1\) and \(\tau=3\).  One exact NVIDIA feature
artifact is generated for 2000 inner-training, 500 validation, and 500 test
windows.  Two Mat\'ern readouts are selected with
\cref{eq:empirical-krr-grid}, refitted on the 2500 non-test windows, and then
frozen.  The hardware subset consists of eight equally spaced indices from the
500-window test block; those indices are stored before any hardware job is
submitted.

The CSMoM estimator uses \(S=192\) local-Pauli snapshots and six deterministic
MoM blocks.  The shadow-basis and reset-trajectory arrays are created from
independent named streams and reused by the local circuit and hardware layers.
Measurement outcomes are independent across executions.  The feature scaler
and readout are always those fitted from exact features.

\subsection{Comparison ladder}
\label{app:hardware-ladder}

For the same eight windows and program, E9 computes:
\begin{enumerate}
  \item \textbf{Aer exact}: analytic expectations from the independent
        \(2n+1\)-qubit density-matrix dilation; this is the mathematical
        feature reference for the subset;
  \item \textbf{Aer shadow-only}: CSMoM samples drawn from the exact Aer channel,
        isolating local-Pauli finite-shadow error;
  \item \textbf{IBM-local trajectories}: the IBM circuit compiler and
        pre-sampled last-reset trajectories executed on an ideal local Aer
        sampler, combining trajectory and shadow sampling without device noise;
        and
  \item \textbf{IBM hardware}: the identical structural plan transpiled and
        executed through IBM Runtime.
\end{enumerate}
The distinction is semantic: Aer exact is exact mixed-channel evolution; Aer
shadow-only estimates that exact state; IBM-local and hardware use one
pre-sampled reset trajectory per shadow snapshot.  No simulator-only noise
instruction is described as a hardware run.

\subsection{Backend and physical-layout policy}
\label{app:hardware-selection}

At the first hardware period, the driver lists operational non-simulator IBM
backends with at least two qubits.  For every connected physical pair
\((q_0,q_1)\), it computes
\[
 s(q_0,q_1)
 =e_{2\mathrm q}(q_0,q_1)
  +\frac{0.25}{2}\bigl(e_{\mathrm{ro}}(q_0)+e_{\mathrm{ro}}(q_1)\bigr),
\]
where unavailable calibration values make the pair ineligible.  The pair with
smallest score is selected; scores within 5\% are broken by backend pending-job
count and then lexicographically by backend and qubit IDs.  Backend-properties
payload, selection candidates, score, and selection time are persisted.

Transpilation uses optimization level one, the named transpiler seed, and the
selected pair as a fixed initial layout.  No post-result remapping is allowed.
The same backend and physical pair are reused in a second period at least 24
hours later when operational.  If they are unavailable, the replacement is
selected by the same rule and is reported as a separate operating point rather
than pooled with the first.

\subsection{Circuit grouping, Runtime jobs, and cost bound}
\label{app:hardware-jobs}

Each snapshot is represented by its window, reservoir, sampled last-reset
suffix, and local measurement basis.  Snapshots with identical tuples share one
circuit and are executed as repeated shots.  With \(N_{\mathrm{HW}}=8\),
\(w=4\), and \(n=2\), at most
\[
 N_{\mathrm{HW}}(w+1)3^n=360
\]
unique groups are possible, while the total primary shots remain
\(N_{\mathrm{HW}}S=1536\) per period.  Ordered groups are split into chunks of
at most 100 Runtime PUBs, giving at most four jobs per period.  Each job handle,
ordered group range, and expected snapshot count is written before waiting for
results, so decoding is resumable.

A failed chunk is retried once with the same circuits and structural plan.
There is no imputation.  A period contributes feature and prediction aggregates
only if all 1536 expected snapshots decode successfully.  Partial periods remain
in the job-success table with their completed counts.  The primary design uses
at most 3072 QPU shots over two periods; the single-retry contingency raises the
hard bound to 4608 shots.  This bound excludes local simulator and transpilation
work.

\subsection{Mitigation and calibration provenance}
\label{app:hardware-mitigation}

The primary E9 protocol uses no measurement mitigation because the current
hardware adapter does not expose a predeclared mitigation implementation.
Mitigation may be added only as a versioned secondary experiment implemented
and tested before any mitigated hardware result is inspected.  It cannot
replace the raw-hardware result.

For every transpiled circuit and job, the artifact records backend name,
backend-properties update time, submission and completion times, physical
layout, depth, two-qubit gate count, total operation counts, Runtime job ID,
shots, expected and decoded snapshots, and terminal status.  Calibration
payloads are stored as raw JSON-compatible snapshots where provider policy
permits, with a checksum otherwise.

\subsection{Metrics, uncertainty, and displays}
\label{app:hardware-metrics}

For every non-exact ladder layer, we report coordinate feature RMSE and the
per-window \(e_{2,i}\) relative to Aer exact.  Hardware is additionally compared
with the matched IBM-local CSMoM realization, which shares bases and
trajectories but not outcomes.  The frozen-readout report includes mean absolute
prediction perturbation, MSE, NRMSE, and relative degradation for both delays.
Uncertainty is summarized by a snapshot bootstrap that resamples complete
snapshot indices within each window and reconstructs CSMoM with six blocks;
the two calibration periods are displayed separately rather than treated as
i.i.d. replicates.  Per-window paired differences are shown alongside the
bootstrap interval.

The main paper receives one compact two-panel figure: feature error across the
four ladder layers and exact-versus-hardware frozen-readout predictions for the
eight windows.  A small adjacent table reports backend, period, layout, median
and maximum depth, total two-qubit gates, circuit groups, shots, and completion
status.  The appendix retains the full per-observable bias, per-window errors,
transpilation rows, calibration metadata, and job ledger.  A null or degraded
result is scientifically valid: it identifies the hardware operating point at
which exact-feature transfer fails.
